\chapter{Аналитическая часть}

В данном разделе будет рассмотрена постановка задачи, особенности анализа PDF-документов,
а также теоретические основы применения регулярных выражений и больших языковых моделей (БЯМ)
для разработки программного обеспечения.

\section{Особенности анализа PDF-документов}

Формат PDF ориентирован на точное отображение документа при печати и просмотре,
поэтому он не хранит логическую структуру в привычном виде (разделы, заголовки, оглавление),
а представляет документ как набор страниц с текстовыми и графическими блоками.
Это осложняет автоматическую проверку структуры: один и тот же визуальный результат
может соответствовать различным внутренним представлениям.

Для решения задачи используется библиотека PyPDF2, позволяющая извлекать текст постранично.
После извлечения текста анализ выполняется на уровне строк, что позволяет ввести упрощённые координаты
обнаруженных фрагментов в виде пары «номер страницы --- номер строки».
Такой подход соответствует требованиям задания и не требует вычисления точных геометрических координат
текстовых блоков.

\section{Условия ошибки в PDF-документах}

Согласно варианту задания, документ считается содержащим ошибку в следующих случаях:
\begin{itemize}
    \item структурный элемент отсутствует в основном тексте документа;
    \item структурный элемент отсутствует в содержании (оглавлении);
    \item используется ошибочное наименование структурного элемента
    (например, «Введение» вместо «ВВЕДЕНИЕ», «Литература» вместо
    «СПИСОК ИСПОЛЬЗОВАННЫХ ИСТОЧНИКОВ», «Библиография», «Источники» и т.п.);
\end{itemize}

\section{Применение регулярных выражений}

Для автоматизации проверки используются регулярные выражения.
Регулярные выражения являются формальным инструментом для описания шаблонов строк и поиска совпадений в тексте.
Их преимущество заключается в детерминированности результата: при фиксированных правилах поиска
программа выдаёт воспроизводимый итог без неоднозначной интерпретации.
В то же время при анализе PDF необходимо учитывать, что извлечённый текст может содержать лишние пробелы,
переносы, разрывы слов и другие артефакты конвертации. Поэтому при разработке шаблонов поиска
закладывается устойчивость к вариативности регистра, пробельных символов и возможных форматных искажений.

\section{Использование больших языковых моделей (БЯМ)}

В рамках выполнения лабораторной работы применялись большие языковые модели (БЯМ).
БЯМ представляют собой нейросетевые модели, обученные на больших корпусах текстов,
которые способны генерировать связный текст, объяснять решения и предлагать программный код по описанию задачи.
Практическая ценность БЯМ в данной работе заключается в ускорении этапов проектирования:
модели помогают формулировать требования, предлагать архитектуру решения, генерировать черновые реализации
и указывать возможные крайние случаи.

\section{Краткое описание используемых БЯМ}

В рамках разработки решения были использованы три большие языковые модели:

\begin{itemize}
    \item \textbf{ChatGPT (GPT-5.2)} — модель общего назначения, которая используется для уточнения требований,
    построения структуры решения и получения развёрнутых объяснений по алгоритму.
    Эта модель эффективно помогает на этапе генерации общей структуры программы, а также в уточнении задач и требований.
    \item \textbf{DeepSeek DeepThink} — модель, ориентированная на последовательное рассуждение и
    проработку логики проверки, включая формулирование правил обнаружения ошибок.
    Она используется для создания устойчивых алгоритмов поиска структурных элементов в тексте, а также для выявления ошибок.
    \item \textbf{Qwen Coder} — модель, специализирующаяся на генерации программного кода и типовой инженерной структуре
    (разделение ядра алгоритма, интерфейса командной строки и модулей вывода результатов).
    Модель Qwen Coder полезна на этапе реализации алгоритма, когда требуется быстро сгенерировать программный код,
    который затем можно доработать и адаптировать под конкретные требования.
\end{itemize}

\section{Вывод}

В этом разделе были рассмотрены особенности анализа PDF-документов и выбран подход к определению координат
обнаруженных фрагментов в формате «номер страницы --- номер строки».
Были сформулированы условия, при которых документ считается содержащим ошибку согласно варианту,
а также обосновано применение регулярных выражений для поиска структурных элементов.
